% =====================================================================
%   3. AgentBench 的构成 (Composition of AgentBench: A Brief Look)
% =====================================================================

\section{AgentBench 的构成：概览}

本节将简要介绍构成 AgentBench 的数据集与环境。
相较于既有的智能体评测基准 \citep{cote2019,fan2022}，AgentBench 专注于
通过思维链（Chain-of-Thought, CoT）\citep{wei2022b,yao2023b} 提示
对 LLM 进行实用性评估，覆盖代码接地、游戏接地与网页接地三大场景。
它们指向 LLM 在自主完成复杂任务上的可期应用方向，并凭借场景多样性，
避免了在 AgentBench 上专用于某一类任务（如代码类）的 LLM 表现出不公平的优势。
受版面所限，本节仅作概念性概述；数据集的构造细节、评测流程与提示词示例，
请参见附录。

\subsection{代码接地类环境}

LLM 具备生成高质量代码 \citep{chen2021} 的能力，因此协助人机交互操作系统
（Operating System）便成为 LLM 智能体的一项重要且务实的使命。
AgentBench 选择了三项任务作为代码与推理类能力的代表性环境。

\paragraph{操作系统（Operating System, OS）}
允许 LLM 在终端中访问并操纵操作系统，是一项既迷人又充满挑战的任务。
尽管将自然语言翻译为 Shell 命令的工作早已有之 \citep{lin2018}，但
几乎少有工作对模型在可执行环境中的表现进行评估。我们的目标，是让 LLM 在
真实可交互的 bash 环境（即 Ubuntu Docker 容器 \citep{merkel2014}）中，
回答具有确定性答案的问题（例如：系统中具有/不具有 home 目录的用户数），
或完成一系列为达成实际目标的连续操作（例如：递归地将所有目录下属于他人
的文件设为只读，而不动我自己创建的文件）。本环境以成功率（Success Rate, SR）
作为评测指标。（详见附录~\ref{app:os}）

\paragraph{数据库（Database, DB）}
数据库分析在日常事务中至关重要却也充满挑战，因而对 LLM 在真实数据库上
通过 SQL 操作的能力进行测评尤为关键。既有研究多侧重于单项流程的局部环节，
例如 SQL 与自然语言之间的翻译 \citep{zhong2017,gao2023,pourreza2023,ruan2023}，
或基于单张小表的问题解答 \citep{nan2021,iyyer2017}，少有工作将完整流水线
作为一个整体加以评估。因此，AgentBench 在真实的 SQL 接口、数据库与
多样化查询上对 LLM 进行评测，力图还原真实世界中的使用场景。
本环境以 SR 作为主要评测指标。（详见附录~\ref{app:db}）

\paragraph{知识图谱（Knowledge Graph, KG）\citep{anonymous2023kg}}
当代知识图谱往往规模庞大（例如 F REEBASE \citep{bollacker2008}
包含超过 4500 万实体与 30 亿事实），这要求智能体具备广泛的能力
\citep{gu2023discriminate}。在此类仅部分可观察的环境中运作，智能体必须在
信息不完备的情形下作出决策，并妥善应对各种内生不确定性——这要求其具备
语言理解（如措辞的细微差别）、任务规划（如将复杂指令拆解为可执行子步骤），
以及工具使用（如与 KG 接口交互）等多种技能。为此，我们提出 KG 作为
评估 AI 智能体决策能力的代表性测试场。本环境以问答（Question Answering）
作为基本任务形式，采用答案 F1 作为评测指标。（详见附录~\ref{app:kg}）

\subsection{游戏接地类环境}

玩电子游戏通常要求智能体具备较强的策略设计、指令遵循与推理能力。
不同于代码类任务，游戏类环境不要求编程专长，而更侧重对常识与世界知识的
全面理解。

\paragraph{数字卡牌游戏（Digital Card Game, DCG）}
游戏通常可作为智能体开发的仿真环境 \citep{hoover2020}。DCG 是面向
纯文本 LLM 的理想选择：它通常包含丰富的卡牌文本描述，要求玩家进行深思熟虑
的策略组合方可取胜，从而考察模型对游戏规则、操作逻辑与基于当前态势的
策略决断能力的理解。在 AgentBench 中，我们采用 2021 年清华大学智能体竞赛
（THUAC）中的简化版 DCG 系统 Aquawar\footnote{\url{https://www.saiblo.net/}}
对 LLM 智能体进行评估。在 Aquawar 中，智能体扮演一方玩家，操控一支由不同
天赋的水族宠物组成的队伍，与另一支算法控制的队伍进行回合制对战。
本环境以 LLM 的胜率作为评测指标。（详见附录~\ref{app:dcg}）

\paragraph{横向思维谜题（Lateral Thinking Puzzles, LTP）}
横向思维谜题（Lateral Thinking Puzzle，亦称情境谜题或俗称"海龟汤"）
\citep{sloane1992,debono1970} 是一项在全球范围内广受欢迎的群体游戏，
鼓励参与者从非常规视角破解谜题。通常，由一人担任主持人，其他人通过
策略性提问来推测谜底，回答限于"是"、"否"或"无关"。在本数据集中，
我们首先构建了一套 LTP 主持人系统以实现自动评判，并汇集了难度各异的
网页谜题，将剧情简化为若干关键节点（即游戏进度）。通过此项评估，我们
期望深入了解 LLM 的横向思维能力深度与敏捷度。
（详见附录~\ref{app:ltp}）

\paragraph{家务整理（House Holding, HH；ALFWorld \citep{shridhar2020b}）}
家务整理等具身游戏环境要求较强的常识接地能力，已成为语言智能体评测的经典
任务 \citep{cote2019}。AgentBench 沿用了经典 ALFWorld
\citep{shridhar2020b}（由成熟的文本游戏工具包 TextWorld \citep{cote2019}
衍生而来）对模型的物理家务环境任务完成能力进行评估。智能体需要完成诸如
"将一只平底锅放到餐桌上"之类的家务任务。我们以 SR 作为评测指标。
（详见附录~\ref{app:hh}）

\subsection{网页接地类环境}

网页已成为人们与现实世界交互的主要界面。因此，评估 LLM 智能体在复杂
网页环境中的行为，对于未来智能体的研发具有重要的现实意义。
本节中，我们改造了两个现有的网页浏览数据集，以对 LLM 进行实用性评测。

\paragraph{网页购物（Web Shopping, WS；WebShop \citep{yao2022}）}
在线购物是现代生活中极为实际且重要的环节。其完整流程——在真实电商网站上
进行检索、浏览并选择心仪商品——要求自主智能体具备强大的推理与决策能力。
WebShop \citep{yao2022} 即是这样一项用于评估语言智能体的在线购物仿真环境。
在其原始工作中评估的是专门训练的模型，而我们则采用纯提示（prompt）方式对
LLM 进行评测。（详见附录~\ref{app:ws}）

\paragraph{网页浏览（Web Browsing, WB；Mind2Web \citep{deng2023}）}
通用网页环境是训练与评估智能体的理想沙盒。Mind2Web \citep{deng2023}
是一项近期发布的通用网页智能体基准，旨在评估能够在不同网站领域内完成复杂
任务的智能体能力，并给定高层级用户指令。它为网页交互设计了可行的动作
（如点击、选择、键入），以支持对 LLM 作为网页智能体的整体评估。
相较于 Mind2Web 原设定，我们进行了必要的适配，使得无需任何额外微调
即可对基于提示词（prompt）的 LLM 进行评估。（详见附录~\ref{app:wb}）
